\section{The whole space}\label{sec:whole}
We now work on $\R^3$, with the data classes $\mathcal X_{\R}$ and
$\mathcal F_{\R}$ from Section~\ref{sec:preliminaries}. The construction
is again spatially local. The new issues are the behavior of negative
Sobolev norms near frequency zero and the absence of a spectral gap.
\subsection{Statement of the density theorem}
\begin{theorem}[Two whole-space Sobolev thresholds]\label{thm:Rmain}
Fix $\nu,T>0$, let $q\in\{1,2\}$, and set $s_q=2/q-3/2$. In the relative $L^q(0,\infty;H^s(\R^3))$ topology on $\mathcal F_{\R}$:
\begin{enumerate}[label=(\roman*),nosep]
\item For every fixed $a\in\mathcal X_{\R}$, $\mathcal B^\R_{\nu,a,T}$ is dense if $s<s_q$.
\item For zero initial velocity, $\mathcal B^\R_{\nu,0,T}$ is dense if and only if $s<s_q$.
\end{enumerate}
Thus the thresholds are $1/2$ for $L^1_tH^s_x$ and $-1/2$ for $L^2_tH^s_x$. Around every given solution smooth through $T$, the approximating solution can have the same initial velocity, the same earlier history, singularity exactly at $T$, and a velocity difference tending to zero in $E_T$.
\end{theorem}

The force classes remain smooth, equipped with the relative norm topology. For $q=2$, the converse uses a regular neighborhood whose radius depends on $T$.

\subsection{Local construction and negative Sobolev estimates}
The compact solution of Theorem~\ref{thm:packet} is already defined on
$\R^3$. Use the scaling~\eqref{eq:scaling} directly, without
periodization. The energy identities~\eqref{eq:packetEscale} and mixed
force identity~\eqref{eq:packetFscale} follow from the same changes of
variables. The vector potential in Lemma~\ref{lem:potential} is local,
so its proof applies in any Euclidean ball. In
Lemma~\ref{lem:correction}, the support, derivative, energy, and mixed
Lebesgue bounds are also local; only the final transfer of fractional
norms to the torus is unnecessary. Thus the same $w_\eps$ and $H_\eps$
are available on $\R^3$, with the same uniform bounds. We record the
result and prove the additional negative-order estimates.
\Needspace{16\baselineskip}
\begin{theorem}\label{thm:Rinsert}
Let $(v,\pi,g)$ be the solution for $a\in\mathcal X_{\R}$ and $g\in\mathcal F_{\R}$, smooth through $T+\delta$ for some $\delta>0$. Fix any nonempty open ball $B\subset\R^3$. For all sufficiently small $\eps>0$, there are $g_\eps\in\mathcal F_{\R}$ and a solution $u_\eps$ such that
\[
 T^\nu_{\max,\R}(a,g_\eps)=T,\quad
 \limsup_{t\uparrow T}\norm{u_\eps(t)}_\infty=\infty,\quad
 u_\eps=v\ \ (0\le t\le T-2\eps^2).
\]
The velocity difference is supported inside $B$ at every $t<T$, and $g_\eps-g\in C_c^\infty(B\times(0,\infty))$. Moreover,
\begin{equation}\label{eq:REclose}
 \norm{u_\eps-v}_{E_T}\le(M+D)\eps^{1/2}+C\eps^{3/2}.
\end{equation}
For $q\in\{1,2\}$ the force difference tends to zero in $L^q_tH^s_x$ whenever $s<2/q-3/2$. All of these conclusions hold for the same family of glued solutions.
\end{theorem}
\begin{proof}
Place the fields in \eqref{eq:scaling} inside $B$, and choose the compact cutoffs of Lemma~\ref{lem:potential}. Define $A,w_\eps,H_\eps$ by \eqref{eq:potential}, \eqref{eq:cutoff} and \eqref{eq:H}. These formulas only use $v$ on a fixed compact neighborhood of $B$ and of time $T$, so the preceding lemmas apply. In particular $v+w_\eps=0$ on a neighborhood of the active building block support. Put
\[
 u_\eps=v+w_\eps+U_\eps,\qquad
 p_\eps=\pi+P_\eps,\qquad
 g_\eps=g+H_\eps+F_\eps.
\]
Writing $b_\eps=v+w_\eps$, expansion of the equation leaves only the cross-advection terms $(b_\eps\cdot\nabla)U_\eps+(U_\eps\cdot\nabla)b_\eps$. They vanish on a neighborhood of the building block support because $b_\eps=0$ there. Outside that support, the building block and its derivatives vanish; smoothness gives the same conclusion at the support boundary. When the building block is inactive it is identically zero. Thus the two terms vanish pointwise everywhere, and the equation is exact. The pressure difference may be chosen to be the compact scalar $P_\eps$: the equation gives an $L^2$ pressure gradient, and applying $I-\PP$ recovers \eqref{eq:Rpressure}. Both force corrections are globally smooth and spacetime compact, including across $T$, so their sum preserves membership in $\mathcal F_{\R}$.

The initial value and earlier history are unchanged. At each $t<T$ the compact blowup solution and cutoff correction lie in $H^\infty$, and the maximum norm of $U_\eps$ is unbounded at $T$. Proposition~\ref{prop:local} identifies the solution with the unique maximal solution. An extension through $T$ would be bounded in $C_tH^2$ on a neighborhood of $T$, hence bounded in $L^\infty_x$ by \eqref{eq:Rproduct}, contradicting the blowup of $U_\eps$. Thus its maximal lifespan is exactly $T$. Equation~\eqref{eq:packetEscale} and the correction estimate in Lemma~\ref{lem:correction} give \eqref{eq:REclose} by the triangle inequality.

It remains to justify every asserted force topology, including negative orders. Let
\[
 \beta(q,s)=\frac2q-\frac32-s.
\]
For $0\le s\le1$, Euclidean scaling and
$\norm{z}_{H^s}\le C_s(\norm{z}_2+\norm{z}_{\dot H^s})$ give
\begin{align}
 \norm{F_\eps}_{L^q_tH^s_x}
 &\le C_{q,s}\bigl(\eps^{2/q-3/2}+\eps^{\beta(q,s)}\bigr),\nonumber\\
 \norm{H_\eps}_{L^q_tH^s_x}
 &\le C_{q,s}\bigl(\eps^{2/q-1/2}+\eps^{\beta(q,s)+1}\bigr).
 \label{eq:RpositiveScale}
\end{align}
The second line uses the uniformly smooth compact rescaled profile of Lemma~\ref{lem:correction}, whose amplitude is $\eps^{-2}$.

For $-3/2<s<0$, every smooth compact profile has finite $\dot H^s$ norm. Indeed, on $|\xi|<1$ its Fourier transform is bounded by its $L^1$ norm and $\int_{|\xi|<1}|\xi|^{2s}\dd\xi<\infty$; on $|\xi|\ge1$ its $L^2$ norm suffices. These bounds are uniform for the rescaled correction profiles, which have common compact support and uniform derivatives. Since
$(1+|\xi|^2)^s\le|\xi|^{2s}$, scaling gives
\begin{equation}\label{eq:RnegativeScale}
 \norm{F_\eps}_{L^q_tH^s_x}\le C_{q,s}\eps^{\beta(q,s)},\qquad
 \norm{H_\eps}_{L^q_tH^s_x}\le C_{q,s}\eps^{\beta(q,s)+1}.
\end{equation}
The time norms run over the entire positive axis. Each rescaled force has time support of length $O(\eps^2)$, including any part after $T$.

For $q=1$, \eqref{eq:RpositiveScale} proves convergence for $0\le s<1/2$; all $s<0$ follow from $\norm{z}_{H^s}\le\norm{z}_2$. For $q=2$, \eqref{eq:RnegativeScale} proves convergence for $-3/2<s<-1/2$. If $s\le-3/2$, choose $r\in(-3/2,-1/2)$ with $r>s$ and use $\norm{z}_{H^s}\le\norm{z}_{H^r}$. This last step avoids making any false homogeneous scaling assertion at indices where a generic compact profile can have an infinite homogeneous norm.
\end{proof}

\subsection{The critical estimate for time-integrable forcing}
\begin{proposition}\label{prop:Rcritical1}
There is a universal $c>0$ such that, for $a\in\mathcal X_{\R}$ and $f\in\mathcal F_{\R}$,
\[
 \norm{a}_{\dot H^{1/2}}+\norm{f}_{L^1_t\dot H^{1/2}_x}<c\nu
 \quad\Longrightarrow\quad T^\nu_{\max,\R}(a,f)=\infty.
\]
In particular, for $a=0$, the ball $\norm{f}_{L^1_tH^{1/2}_x}<c\nu$ consists of globally regular inputs.
\end{proposition}
\begin{proof}
This is the $p=r=2$ case of Danchin's forced critical small-data theorem
\cite[Theorem~2.3.1, pp.~47--49]{danchin}. Indeed,
$\dot B^{1/2}_{2,2}=\dot H^{1/2}$ and
$\|f\|_{\widetilde L^1\dot H^{1/2}}\le\|f\|_{L^1\dot H^{1/2}}$;
the Leray projection is bounded on these spaces. The theorem gives a global
critical solution, and persistence of regularity for the smooth data in our
classes identifies it with the maximal classical solution of
Proposition~\ref{prop:local}. Finally
$\|f\|_{\dot H^{1/2}}\le\|f\|_{H^{1/2}}$ gives the stated open ball.
\end{proof}

\subsection{The critical estimate for square-integrable forcing}
\begin{proposition}\label{prop:Rcritical2}
For each $\nu,S>0$ there is $r_{\nu,S}>0$ such that
\[
 f\in\mathcal F_{\R},\qquad
 \norm{f}_{L^2(0,\infty;H^{-1/2})}<r_{\nu,S}
 \quad\Longrightarrow\quad T^\nu_{\max,\R}(0,f)>S.
\]
One can take $r_{\nu,S}=c\nu^{3/2}e^{-C\nu S}$ for suitable universal positive constants $c,C$.
\end{proposition}
\begin{proof}
Use the inhomogeneous heat estimate
\cite[Theorem~2.2.5, p.~44]{danchin} and the critical bilinear fixed-point
argument \cite[Lemma~2.3.2 and the proof of Theorem~2.3.1, pp.~48--49]{danchin}.
We spell out the low-frequency and viscosity adjustments. Normalize viscosity
by $u(t,x)=\nu v(\nu t,x)$ and $F(\tau,x)=\nu^{-2}f(\tau/\nu,x)$.
On $0\le\tau\le L:=\nu S$, set
\[
 X_L=C([0,L];H^{1/2})\cap L^2(0,L;H^{3/2}),
 \qquad
 \mathcal H F(t)=\int_0^t e^{(t-r)\Delta}\mathbb P F(r)\dd r.
\]
The cited heat estimate and the Sobolev product law give, with the sum norm
on $X_L$,
\[
 \|\mathcal H F\|_{X_L}\le Ce^{CL}\|F\|_{L^2H^{-1/2}},\qquad
 \|\mathcal B(v,w)\|_{X_L}\le Ce^{CL}\|v\|_{X_L}\|w\|_{X_L},
\]
where $\mathcal B(v,w)=-\mathcal H\operatorname{div}(v\otimes w)$.
For the second bound use $X_L\hookrightarrow L^4_tH^1_x$ and
$H^1\cdot H^1\hookrightarrow H^{1/2}$; the exponential is a convenient
upper bound for the finite-time constants in the inhomogeneous heat estimate.
Thus the cited contraction lemma applies if
$\|F\|_{L^2H^{-1/2}}<ce^{-CL}$, after changing universal constants.
Since
$\|F\|_{L^2(0,L;H^{-1/2})}=\nu^{-3/2}\|f\|_{L^2(0,S;H^{-1/2})}$,
the stated radius suffices. Persistence of smoothness and local existence
at the final time give lifespan strictly greater than $S$.
The finite-time inhomogeneous estimate is essential: small
$H^{-1/2}$ norm does not bound the homogeneous negative norm at low frequencies.
\end{proof}

\begin{proof}[Proof of Theorem~\ref{thm:Rmain}]
Fix $a,g$ and a positive radius in the indicated force norm. There are two cases. If $T^\nu_{\max,\R}(a,g)\le T$, take $f=g$. Otherwise choose $\delta>0$ so that the given smooth solution exists through $T+\delta$, and apply Theorem~\ref{thm:Rinsert}. For $s<s_q$, its force difference is smaller than the prescribed radius when $\eps$ is sufficiently small, its initial velocity is exactly $a$, and its lifespan is exactly $T$. This proves density for each fixed smooth initial velocity.

For zero initial velocity and $q=1$, Proposition~\ref{prop:Rcritical1} provides an open ball about zero disjoint from $\mathcal B^\R_{\nu,0,T}$ in $L^1_tH^{1/2}_x$. The inclusion $H^s\hookrightarrow H^{1/2}$ with norm at most one for $s\ge1/2$ gives the same obstruction in each stronger metric. For $q=2$, use Proposition~\ref{prop:Rcritical2} with $S=T$ and then $H^s\hookrightarrow H^{-1/2}$ for $s\ge-1/2$. These nonempty relative open balls prove non-density. The energy and earlier-history assertions are part of the gluing theorem.
\end{proof}

\subsection{Compact support, rapid decay, and completed force spaces}
Define the following two subclasses of $\mathcal F_{\R}$:
\begin{align*}
 \mathcal F_c&=C_c^\infty(\R^3\times(0,\infty);\R^3),\\
 \mathcal F_{\mathrm{rd}}
 &=\left\{f\in C^\infty(\R^3\times[0,\infty)):
 \sup_{x\in\R^3,\ t\ge0}(1+|x|+t)^N
       |\partial_x^\alpha\partial_t^jf(x,t)|<\infty
 \text{ for all }N,\alpha,j\right\}.
\end{align*}
Vector values are understood in the second line. The initial class corresponding to the rapid-decay formulation is $\mathcal S_\sigma=\mathcal S(\R^3;\R^3)\cap L^2_\sigma$. These are exactly the decay requirements relevant to \cite[equations (4)--(5)]{clay}; a common numerical bound on their seminorms is not imposed.

\begin{corollary}\label{cor:Rclasses}
Theorem~\ref{thm:Rmain} remains valid with $\mathcal F_{\R}$ replaced by $\mathcal F_c$ or $\mathcal F_{\mathrm{rd}}$. In particular it holds for every fixed $a\in\mathcal S_\sigma$ in the rapid-decay class, with the complete if-and-only-if classification at $a=0$.
\end{corollary}
\begin{proof}
The initial value is preserved exactly, and the force correction in Theorem~\ref{thm:Rinsert} is spacetime compact. Addition of that correction preserves each stated force class, including all rapid-decay bounds. The two-case density proof therefore stays within the chosen subclass. Each critical regular ball contains zero and remains a nonempty relative open ball in that subclass. The local framework only requires $a\in\mathcal X_{\R}$, so it applies in particular to Schwartz initial data. Nothing here requires the nonzero given velocity or its pressure to be spatially compact.
\end{proof}

The force spaces in the next proposition are natural for the energy
class. With $V=H^1(\R^3;\R^3)\cap L^2_\sigma$, a finite-energy velocity
belongs to $L^\infty_tL^2_\sigma\cap L^2_tV$. A full distribution in
$H^{-1}$ restricts to an element of $V'$; the latter records only its
action on divergence-free tests. This distinction is discussed in
\cite[Definition~1 and Remark~4]{berselli}. Our realization of
$\dot H^{-1}$ gives the pairing
\[
 |\ip{h}{z}|\le\norm{h}_{\dot H^{-1}}\norm{\nabla z}_2,
 \qquad
 \left|\int_0^T\ip{f}{u}\dd t\right|
 \le\norm{f}_{L^2_t\dot H^{-1}_x}\norm{\nabla u}_{L^2_{t,x}}.
\]
These inequalities follow by Fourier Cauchy--Schwarz, first for Schwartz
functions and then by completion. The inhomogeneous $H^{-1}$ norm instead
pairs with $H^1$, while $L^1_tL^2_x$ pairs with $L^\infty_tL^2_x$.

\begin{proposition}\label{prop:Renergy}
Fix $a\in\mathcal X_{\R}$ and $\nu,T>0$. Smooth compact forces with $T^\nu_{\max,\R}(a,f)\le T$ are dense in each full Bochner space $L^q(0,\infty;H^s(\R^3))$ for $q\in\{1,2\}$ and $s<s_q$. They are also dense in $L^2(0,\infty;\dot H^{-1}(\R^3))$.

For every given smooth solution in Theorem~\ref{thm:Rinsert}, one may simultaneously arrange
\[
 \norm{u_\eps-v}_{E_T}\longrightarrow0,\qquad
 \norm{g_\eps-g}_{L^1_tL^2_x}
 +\norm{g_\eps-g}_{L^2_tH^{-1}_x}
 +\norm{g_\eps-g}_{L^2_t\dot H^{-1}_x}\longrightarrow0.
\]
The homogeneous norm here is a norm of the compact difference; the given force itself need not belong to that homogeneous force space.
\end{proposition}
\begin{proof}
The standard density of Schwartz functions in $H^s$
\cite[Chapter~4, Section~1, Exercise~1]{taylor2011}, followed by smooth spatial
cutoffs, gives density of $C_c^\infty$ for every real $s$. For the homogeneous
space, approximate $|\xi|^{-1}\widehat h$ by smooth annular Fourier data and
then use the same cutoffs. The only additional low-frequency estimate is
\begin{equation}\label{eq:Rnegative-cutoff}
 \norm{k}_{\dot H^{-1}}^2\le
 C\norm{k}_1^2\int_{|\xi|<1}|\xi|^{-2}\dd\xi+\norm{k}_2^2
 \le C'\norm{k}_1^2+\norm{k}_2^2.
\end{equation}
Applied to the Schwartz cutoff tails, this proves spatial compact-smooth
density in $\dot H^{-1}$. The constructions preserve real vector fields.
For time dependence, Hunter~\cite[Proposition~6.29]{hunterpde} gives density
of finite sums $\sum_j\varphi_j(t)b_j$ with $\varphi_j\in C_c^\infty$ in
$L^q(I;X)$, $q<\infty$. Truncate $(0,\infty)$ to $[1/N,N]$ and approximate
each $b_j$ by a spatially compact smooth field. The resulting finite sums
are jointly smooth and compactly supported in space and positive time.

For the inhomogeneous density of forces producing breakdown, first choose such a smooth compact force within half the prescribed radius of the target. Corollary~\ref{cor:Rclasses} supplies a smooth compact force producing breakdown within the other half. This is density of a smooth subset in the completion, not a definition of classical breakdown for every rough force.

For the homogeneous gluing estimate, apply the Fourier scaling calculation of \eqref{eq:RnegativeScale} at $s=-1$. Explicitly, an amplitude $\eps^{-a}$ has spatial $\dot H^{-1}$ factor $\eps^{5/2-a}$, because
$\widehat{\eps^{-a}h(\,\cdot/\eps)}(\xi)=\eps^{3-a}\widehat h(\eps\xi)$.
The time $L^2$ factor is $\eps$. The source force has $a=3$, while the correction has $a=2$ and uniformly bounded compact rescaled profiles; their $\dot H^{-1}$ norms are uniformly finite by \eqref{eq:Rnegative-cutoff}. Hence
\[
 \norm{F_\eps}_{L^2_t\dot H^{-1}_x}\le C\eps^{1/2},\qquad
 \norm{H_\eps}_{L^2_t\dot H^{-1}_x}\le C\eps^{3/2}.
\]
Given a compact smooth force, if its lifespan is at most $T$ it already belongs to the required breakdown set. Otherwise glue the building block and use these estimates. Combining this relative approximation with the proved compact-smooth Bochner density establishes density in $L^2_t\dot H^{-1}_x$. Finally these estimates, \eqref{eq:REclose}, and the $q=1,s=0$ and $q=2,s=-1$ cases of the gluing estimates prove the simultaneous convergence for the same family.
\end{proof}

As on the torus, every $a\in\mathcal X_{\R}$ occurs as the initial velocity of a solution producing breakdown for some smooth force. Thus the projection of the initial data and forces producing breakdown onto $\mathcal X_{\R}$ is the whole space, with quantifier order $\forall a\,\exists f$. The preceding approximation concerns finite-energy whole-space solutions and uses one rescaled copy of $U$ in each perturbation.

\subsection{Cell averages on prescribed grids}\label{sec:grids}
The spaces used for energy estimates should be distinguished from the
additional regularity assumed in numerical error estimates. For example,
on a bounded no-slip computational domain one uses the solenoidal energy
space $H$ with zero normal trace and $V=H^1_0\cap H$; finite-dimensional
velocity and pressure spaces discretize this setting. Guermond, Minev and
Shen~\cite[Section~2 and Theorem~3.2]{gms} state these spaces and the
additional smoothness required for their error estimate. The initial
pressure is constrained by the equation and boundary conditions. These
bounded-domain conventions do not impose boundary conditions at infinity
in the whole-space result below.

To relate the continuous norms to a concrete numerical observation, let $\mathcal T_h$ be a complete uniform Cartesian grid of $\R^3$ with positive mesh widths. For a locally integrable vector field, define the cell-observation map with sequence codomain
\[
 A_h:L^1_{\mathrm{loc}}(\R^3;\R^3)\longrightarrow(\R^3)^{\mathcal T_h},
 \qquad (A_hz)_C=\frac1{|C|}\int_Cz(x)\dd x.
\]
Only coordinatewise equality in this codomain is needed. If a sequence norm is desired for $z\in L^2$, the natural volume-weighted norm obeys
$\sum_{C\in\mathcal T_h}|C|\,|(A_hz)_C|^2\le\norm{z}_2^2$
by Cauchy--Schwarz in each cell and summation. Each grid has infinitely many cells. The next result concerns a finite number of such grids and equality on every cell of each grid, at every time before blowup.

\begin{theorem}\label{thm:Rgrid}
Under the hypotheses on the given smooth solution of Theorem~\ref{thm:Rinsert}, fix a finite family of complete uniform Cartesian grids. The glued solutions may be chosen so that, for every grid in that family,
\[
 A_hu_\eps(t)=A_hv(t),\qquad A_hg_\eps(t)=A_hg(t)
 \qquad(0\le t<T),
\]
while $T^\nu_{\max,\R}(a,g_\eps)=T$ and the energy and force convergences in Proposition~\ref{prop:Renergy} hold. All differences are supported inside a ball contained in one cell of every grid, except for an optional spatially constant pressure gauge.
\end{theorem}
\begin{proof}
The union of faces in one grid is a closed, locally finite union of planes. For the fixed finite family their union is closed and has measure zero. Choose a point outside that union and a ball with closure inside one cell of each grid. Apply Theorem~\ref{thm:Rinsert} inside that ball, using the compact pressure representative. Put $\delta u=u_\eps-v$, $\delta p=p_\eps-\pi=P_\eps$, and $\delta g=g_\eps-g$.

Fix one containing cell $C$. For every $t<T$, $\delta u$ is smooth, divergence free and compactly supported in its interior. For component $j$,
\[
 \delta u_j=\nabla\cdot(x_j\delta u),\qquad
 \int_C\delta u_j\dd x=\int_{\partial C}x_j\delta u\cdot n\dd S=0.
\]
Subtract the two momentum equations and integrate to obtain
\begin{equation}\label{eq:gridforce}
 \int_C\delta g\dd x
 =\frac{\dd}{\dd t}\int_C\delta u\dd x
 +\int_{\partial C}\bigl(u_\eps\otimes u_\eps-v\otimes v
                         +\delta p\,I-\nu\nabla\delta u\bigr)n\dd S.
\end{equation}
The time derivative is zero by the preceding identity. All differences vanish in a neighborhood of the cell boundary, so the surface term is zero. A spatially constant pressure gauge adds a multiple of $\int_{\partial C}n\dd S=0$. The averages therefore agree in $C$; in every other cell they agree by support. The same construction fits the containing cell of each grid, and the convergence estimates are unaffected by this fixed upper bound on $\eps$.
\end{proof}

A deterministic procedure using only the initial velocity cell averages
and the force cell averages for $0\le t<T$ receives identical data in the
two cases. Other observations, such as point values or pressure, may
distinguish them. The perturbation depends on the fixed grid family and
develops larger amplitudes and derivatives as its support shrinks. The
result therefore concerns the information in the prescribed averages,
rather than convergence under refinement for one fixed smooth problem.
